\documentclass[reprint,amsmath,amssymb,aps,prd,twocolumn]{revtex4-2}

\usepackage{graphicx}
\usepackage{dcolumn}
\usepackage{bm}
\usepackage{hyperref}
\usepackage{xcolor}

\begin{document}

\title{Kinematic and Optical Equivalence to Spacetime Curvature:\\A Computational Proof-of-Concept for Flat-Space Relativity}

\author{R. Anastasov}
\affiliation{Anastasov Theory Research Initiative}
\date{\today}

\begin{abstract}
The prevailing astrophysical paradigm universally explains the anomalous orbital precession of Mercury, the deflection of starlight by gravity, the Shapiro time delay, and GPS temporal dissonance mathematically through the metric tensor formalism of General Relativity (curved 4-dimensional spacetime). The explicit goal of this paper is not to invalidate those well-documented astronomical observations, nor to alter empirical data, nor to falsify Einstein's highly successful mathematics. Instead, the purpose of this work is to present a stringent proof-of-concept demonstrating that a functionally identical, overlapping mathematical alternative exists. We assert and demonstrate computationally that the observed relativistic outcomes can be immaculately modeled in a flat Euclidean space, treating gravitational anomalies mechanistically---through physical kinematics, specific angular momentum, and the optical resistance properties of a polarizable vacuum ground state.
\end{abstract}

\maketitle

\section{Philosophical Motivation: What is ``Curved Space'' and ``Curved Time''?}

Physics has long struggled with the epistemological boundary between reality and descriptive geometry. According to purely geometric General Relativity (Einstein’s formalism), gravitational mass structurally twists the invisible fabric of spacetime. Planets follow straight trajectories through warped background dimensions, and ``time itself'' is believed to stretch or compress like a tangible fabric.

However, an alternative physical philosophy existed concurrently, advocated implicitly by thinkers like Lorentz, Poincaré, and Dicke. This philosophy contends a starkly different reality:
\begin{enumerate}
    \item \textbf{Space is flat and unbending}, but acts as a dynamic energy medium (e.g., polarizable vacuum).
    \item \textbf{Time does not curve or stretch.} Time is nothing more than our measurement of the absolute rate of physical processes (e.g., chemical reactions, photon propagation, electron transitions).
\end{enumerate}

When we say ``time slows down'' near a dense black hole, we must assert fundamentally: We firmly agree that every clock made of known matter will truly advance at a different, heavily delayed rate. However, this happens not because an abstract axis of ``Time'' has magically elongated, but because atomic clocks are physical machines. The energetic particles inside those clocks encounter extreme mechanical stress tracking through intense energy gradients or high-speed kinematics. Thus, the physical machinery objectively slows down, giving the false illusion that geometry twisted.

\section{Unifying The Solar System Tests}

\subsection{The Failure of Simple Scalar Upgrades}

Our initial computational attempts sought to upgrade Newtonian physics ($F_{n} = GmM/r^2$) into relativity by applying a basic energetic-momentum multiplier rooted in the total systemic velocity:
\begin{equation}
a_{new} = a_{Newton} \left( 1 + \frac{v^2}{c^2} \right)
\end{equation}
This simplistic scalar multiplier famously fails, yielding roughly $\sim 14''$/cy of precession for Mercury---only $1/3$ of the necessary astronomical rotation of the perihelion. It does not account for the complex asymmetry required to torque an elliptical orbit.

\subsection{The Switch to Specific Vector Angular Momentum}

To correct the model, we observed that gravitational perturbation does not correlate purely with radial kinetic energy (falling straight toward the Sun), but specifically with the \emph{transverse} motion of the particle. The faster the particle cuts \emph{across} the gravitational flux, the stronger the induced anomalous effect.

Mathematically, this transverse state is captured by the Specific Angular Momentum, a localized vector quantity independent of scalar uniform translation:
\begin{equation}
\vec{h} = \vec{r} \times \vec{v}
\end{equation}
This derivation leads us to the precise, corrected acceleration equation acting fundamentally within a flat space grid:
\begin{equation}
\vec{a} = -\frac{GM}{r^3}\vec{r} \cdot \left[ 1 + K(v) \frac{h^2}{c^2 r^2} \right]
\end{equation}

\subsection{Resolving the Clash of Parameters: Light vs. Matter}

Even with the transverse vector equation, physics encounters a critical conflict:
\begin{itemize}
    \item To rotate \textbf{Mercury's} perihelion correctly ($43''$/cy), the vector multiplier \textbf{must be exactly $K = 3.0$}.
    \item To deflect \textbf{starlight} correctly at the limb of the sun ($1.75''$), the vector multiplier \textbf{must be exactly $K = 1.5$}.
\end{itemize}

Academia typically interprets this divergence as proof that two separate dimensional planes (Space vs Time metrics) are operating simultaneously. However, our flat-space simulation introduces the dynamic interpolation factor $K(v)$, which bridges both states seamlessly, attributing the difference mechanically to the intrinsic speed of the observing body $v$:
\begin{equation}
K_{dynamic}(v) = 3.0 - 1.5 \left( \frac{v^2}{c^2} \right)
\end{equation}

\textbf{The Computational Proof:}
\begin{itemize}
    \item \textbf{For Massive Planets ($v \ll c$):} The velocity ratio $v^2/c^2 \approx 0$. The equation resolves cleanly to $K = 3.0$. The simulated orbit yields precisely $43''$/cy.
    \item \textbf{For Light / Radiation ($v = c$):} The ultimate velocity ratio $c^2/c^2 = 1$. The equation resolves cleanly to $K = 3.0 - 1.5 = 1.5$. The simulated photon path traces a perfect $1.75''$ deflection arc.
\end{itemize}

This $K_{dynamic}$ unifies immense stellar bodies and massless photons under a single, non-geometric rule!

\section{The Polarizable Vacuum and Shapiro Delay}

To preserve strict flat-space logic without crashing into super-luminous speeds, we invoked the profound \textbf{Polarizable Vacuum} formalism. A massive celestial body dramatically increases the ``dielectric optical density'' (the refractive index $n$) of the empty spatial vacuum immediately surrounding it:
\begin{equation}
n(r) = 1 + \left(\frac{K_{dynamic}(c)}{1.5}\right) \frac{2GM}{rc^2} = \mathbf{1 + \frac{2GM}{rc^2}}
\end{equation}

The coordinate speed of light locally becomes $v_{light} = c / n(r)$. Photons bound toward the Sun are not accelerated by a fictional geometric dip; their internal wave-front simply degrades mechanically against a denser local vacuum field.
When modeled via differential propagation $(dt/dx = 1/v_{light})$ in our script spanning Earth to Venus, the simulation retrieved identical empirical data: exactly $\mathbf{\sim 115\ \mu s}$ of round-trip radar delay.

\section{Gravitational Redshift: Universal vs. Quantum Components}

In this framework, gravitational redshift is not caused by the curvature of spacetime, but by the change in the local refractive index of the vacuum $n(\varphi)$.

\subsection{The Universal Component ($z_{universal}$)}
If matter is fundamentally composed of field excitations or ``trapped light'' (standing waves/vortices), its internal dynamic frequency $\omega$ is governed by the local speed of light:
\begin{equation}
\omega \propto c_{loc} = \frac{c_0}{n}
\end{equation}
As a body approaches a massive object, $n$ increases, $c_{loc}$ decreases, and all atomic processes slow down identically. This explains why \textbf{all} spectral lines (Hydrogen, Metals, etc.) show a consistent baseline redshift (e.g., 23.9 km/s for 40 Eri B).

\subsection{The Quantum Differential Component ($z_{alpha}$)}
Superimposed on the universal slowing of $c$ is the change in the fine-structure constant $\alpha$, which depends on the vacuum permittivity $\varepsilon_0(\varphi)$:
\begin{equation}
\alpha = \frac{e^2}{4\pi \varepsilon_0 \hbar c}
\end{equation}
Since $\varepsilon_0$ and $c$ vary differently in the presence of a potential $\varphi$, $\alpha$ changes, shifting electronic energy levels according to their sensitivity coefficients ($q$):
\begin{equation}
\Delta \omega_{q} = q \cdot \left[ \left( \frac{\alpha}{\alpha_0} \right)^2 - 1 \right]
\end{equation}
The total observed redshift is:
\begin{equation}
z_{total} = z_{universal} + z_{alpha}(q)
\end{equation}

\section{Grounding Reality: Mechanical Temporal Dilation in GPS Orbits}

The most severe and definitive practical test of ``spacetime curvature'' occurs directly above our heads inside Global Positioning Satellites (GPS). Atomic GPS clocks are physical machines physically altering their tick rate in response to shifting drag potentials.

\textbf{1. Velocity Stress (The Kinematic Slowing):} 
A GPS satellite travels uniformly at $3.87$ km/s. Orbiting forward at immense speed leaves internal energetic events with diminished capacity to tick transversely. The clock's physical mechanisms undergo localized kinematic drag, dropping ticks:
\begin{equation}
\Delta t_{kinematic} \approx -\frac{1}{2} \frac{v_{sat}^2}{c^2} \quad \implies \quad -7.1\ \mu\text{s/day}
\end{equation}

\textbf{2. Vacuum Density Relief (The Gravitational Speedup):} 
As altitude expands, the gravitational dielectric density (the $n$-index of the vacuum) rapidly decreases. Because atomic clocks are fundamentally governed by electromagnetic transitions (Matter as Light), their internal operational speed is dictated by the local speed of light $c_{loc} = c_0 / n(\varphi)$. With a lower $n$ at high altitudes, $c_{loc}$ increases:
\begin{equation}
\Delta t_{gravitational} \approx \frac{GM}{c^2} \left[ \frac{1}{R_{earth}} - \frac{1}{R_{sat}} \right] \implies +45.7\ \mu\text{s/day}
\end{equation}

\textbf{The Simulation Verdict:}
Subtracting the kinematic delay from the density speedup generates a net clock advance of $\mathbf{38.6\ \mu s}$/day. Without accounting for this mechanical behavior, the GPS network would accumulate a map error of nearly $11.6$ kilometers per single day.

\section{Discussion and Conclusion}

We have demonstrated that the core macroscopic predictions of General Relativity---perihelion precession, light deflection, Shapiro delay, and GPS temporal dilation---can be successfully modeled without abandoning flat Euclidean space. By treating the vacuum as a polarizable dielectric medium where $c_{loc} = c_0 / n(\varphi)$, and recognizing matter as fundamentally light-like field configurations, the apparent ``curvature of time and space'' emerges as an optical and kinematic consequence of the variable vacuum potential.

Furthermore, this framework bridges the gap between classical kinematics and quantum scaling. While standard General Relativity demands a strict adherence to the geometric Equivalence Principle, the $\varepsilon_0(\varphi)$ model naturally predicts a spectroscopic violation via the element-specific sensitivity to the fine-structure constant. The empirical identification of $q$-dependent residuals in white dwarf spectra provides a robust, falsifiable baseline for future observational verification.

\end{document}
